\subsection{Rotation and hierarchical codes}

JHQ first applies a square orthogonal transform
$\Pi\in\mathbb R^{d\times d}$ to both database vectors and queries
\cite{jhq}. Write $y=\Pi x$ and $q'=\Pi q$. In real arithmetic,
\begin{equation}
 \begin{aligned}
 \Pi^\top\Pi=I
 &\Longrightarrow
 \|\Pi(q-x)\|^2\\
 &=(q-x)^\top\Pi^\top\Pi(q-x)=\|q-x\|^2.
 \end{aligned}
 \label{eq:2-isometry}
\end{equation}
Thus rotation preserves dimension and distances before quantisation.
Finite-precision matrix multiplication is evaluated separately in Chapter 4.

The primary code gives a coarse reconstruction $\widehat p(x)$ of $y$.
In the evaluated configuration, each eight-coordinate subspace uses one
byte: each coordinate takes one of two levels, $-a$ or $+a$, so the byte
records eight signs. More generally, $M$ subspaces have dimension
$D_s=d/M$, with $B$ bits selecting one of $K=2^B$ combinations of
$L=2^{B/D_s}$ scalar levels. For supported parameters with integer $L$,
JHQ uses the Cartesian product
\begin{equation}
 a_i=\sigma\Phi^{-1}\!\left(\frac{i-1/2}{L}\right),\quad
 i=1,\ldots,L,\qquad
 \mathcal C_m=\{a_1,\ldots,a_L\}^{D_s}.
 \label{eq:2-cartesian}
\end{equation}
Here $\Phi$ is the standard normal cumulative distribution function and
$\sigma$ is a scale estimated from vector energy. At $B=D_s=8$, this
gives $L=2$, $\mathcal C_m=\{-a,+a\}^8$ and
$a\simeq0.67449\sigma$. Unit-normalised inputs motivate a scale near
$1/\sqrt d$. The transform and hierarchical representation are inherited
from JHQ.

Residual coding approximates the remaining error,
$r=y-\widehat p$, by $\widehat r$, giving
$\widehat y=\widehat p+\widehat r$. The evaluated residual quantiser learns
scalar levels shared across coordinates within each subspace and uses one
byte per coordinate. Each vector therefore stores $d/8$ primary bytes
and $d$ residual bytes before metadata. A smaller refinement budget reduces
residual reads, while these stored codes remain unchanged.

This residual is relative to the primary reconstruction. By contrast,
residual-coded IVF-PQ subtracts the centroid $\mu_{\ell(x)}$ of the list
containing $x$:
\[
 r_{\rm IVF\text{-}PQ}=x-\mu_{\ell(x)},\qquad
 r_{\rm JHQ}=\Pi x-\widehat p(x).
\]
JHQ uses IVF for routing; its reconstructed vector does not add an IVF centroid.

\subsection{Primary and refined scoring}

Primary screening compares $q'$ with the coarse reconstruction:
\[
 \primarydist=\|q'-\widehat p\|^2.
\]
Only shortlisted candidates are rescored with both primary and residual
codes. The implemented refined score is
\begin{equation}
 \refineddist=\|q'-\widehat p-\widehat r\|^2+\|q'\|^2.
 \label{eq:2-refined}
\end{equation}
The extra term $\|q'\|^2$ is constant across candidates for one query,
so the score orders reconstructed vectors by squared distance. It need
not reproduce the raw-vector ordering used for GT. The score expansion
appears in \cref{eq:5-refined-constant}.

\par\medskip\noindent\begin{minipage}{\linewidth}
\centering\captionsetup{hypcap=false}
\resizebox{\linewidth}{!}{\begin{tikzpicture}[
 font=\small\linespread{1}\selectfont,
 box/.style={draw=ThesisNavy!65,rounded corners=2pt,align=center,
 text width=2.65cm,minimum height=1.8cm,inner sep=4pt},
 arr/.style={-{Latex[length=2mm]},draw=ThesisNavy!80,thick}]
\node[box] (x) at (0,0) {Database $\mathcal X$\\$N$ vectors};
\node[box] (u) at (3.6,0) {IVF lists\\$U(q)$: $n_c$ vectors};
\node[box,fill=ThesisNavy!7] (c) at (7.2,0) {Primary top-$c$\\$C_\alpha(q)$\\up to $c$ vectors};
\node[box,fill=ThesisNavy!7] (r) at (10.8,0) {Refined top-$k$\\$\mathcal R_\alpha(q)$};
\draw[arr] (x) -- node[above,font=\footnotesize] {IVF} (u);
\draw[arr] (u) -- node[above,font=\footnotesize] {$D_{\rm P}$} (c);
\draw[arr] (c) -- node[above,font=\footnotesize] {$D_{\rm H}$} (r);
\node[align=center,text width=6.4cm] at (2.6,-2.0)
 {Query $q'=\Pi q$ stays FP32.\\Primary codes score the broad scan.};
\node[align=center,text width=5.8cm] at (9.7,-2.0)
 {Residual $r=\Pi x-\widehat p$.\\Survivors decode $\widehat p+\widehat r$.};
\end{tikzpicture}
}
\captionof{figure}[Candidate and scoring cascade]{Candidate and scoring cascade in JHQ.
IVF selects lists, primary scores form a shortlist, and residual codes
refine that shortlist before the final top-$k$ selection.}
\label{fig:2-cascade}
\end{minipage}\par\medskip

With $n_c$ routed vectors, $M$ primary subspaces and up to $c$ survivors,
the score work is
\[
 W_{\rm primary}=O(n_cM)\ \text{lookups},\qquad
 W_{\rm refine}=O(cd)\ \text{coordinates}.
\]
Memory access and candidate selection also contribute to search time.

\subsection{Modelling assumptions}

Random rotation can spread coordinate energy, motivating the Gaussian-inspired
levels above. Distance preservation does not imply independent, zero-mean
Gaussian coordinates: a fixed transform maps a database mean $\mu$ to
$\Pi\mu$, with $\|\Pi\mu\|=\|\mu\|$. Chapter 8 examines centring as an
alternative coding choice.

The primary level is a Gaussian mid-quantile, rather than the level that
minimises reconstruction mean-squared error (MSE). For
$X\sim\mathcal N(0,\sigma^2)$ and a zero decision threshold, the latter is
\[
 \begin{aligned}
 a_{\rm MSE}
 &=\arg\min_a\mathbb E[(X-a)^2\mid X>0]\\
 &=\mathbb E[X\mid X>0]
 =\sigma\sqrt{2/\pi}\simeq0.79788\sigma.
 \end{aligned}
\]
These values serve different objectives. For fixed sign codes, any
positive common primary scale preserves primary ranking
(\cref{eq:4-sign-ranking}); residual reconstruction still depends on the
scale. Thus the algebraic scoring properties do not establish optimal
quantisation or raw-vector recall.
